\section{System F normalization: complete candidate argument}
\label{sec:annex-p07-system-f-normalization}

\begin{annexitem}{annex:p7.system-f.substitution-reduction}
 {thm:p07-system-f-sn-substitution-reduction}{Both substitutions and all reduction positions}
\begin{proof}
First record the syntactic commutations used in the proof. All terms below
are raw, and substitutions are simultaneous. Binders are renamed fresh
before applying an identity. At a term-beta redex, type substitution gives
\[
 ((\lambda x{:}\alpha.P)N)\tau
  =(\lambda x{:}\alpha\tau.P\tau)(N\tau)
  \to_{\mathsf F}(P\tau)[N\tau/x]
  =(P[N/x])\tau.
\]
Term substitution, with \(x\) fresh for its domain and images, gives
\[
 ((\lambda x{:}\alpha.P)N)\sigma
  =(\lambda x{:}\alpha.P\sigma)(N\sigma)
  \to_{\mathsf F}(P\sigma)[N\sigma/x]
  =(P[N/x])\sigma.
\]
At a type-beta redex choose its binder \(a\) fresh for the domains and
free type variables of all other substitution images. Then
\[
 ((\Lambda a.P)[\beta])\tau
  =(\Lambda a.P\tau)[\beta\tau]
  \to_{\mathsf F}(P\tau)[\beta\tau/a]
  =(P[\beta/a])\tau,
\]
and
\[
 ((\Lambda a.P)[\beta])\sigma
  =(\Lambda a.P\sigma)[\beta]
  \to_{\mathsf F}(P\sigma)[\beta/a]
  =(P[\beta/a])\sigma.
\]
The equalities of contracta follow from structural substitution induction.
For a term variable, distinguish the substituted variable from the other
variables, and use the image specified by the relevant substitution.
Term application distributes the equality to both children. A lambda
abstraction distributes it to its annotation and freshly renamed body.
A type abstraction uses a type binder fresh for both substitutions; a type
application distributes it to its term and type argument. In the type
components, the corresponding induction distinguishes a type variable,
an arrow with its two children, and a universal type with a freshly renamed
binder. These cases also prove that substitution acts once on an inserted
image, rather than applying itself again to that image.

For a contraction inside a larger term, induction on its surrounding
context propagates the same genuine contraction through an application's
function or argument, a lambda body, a type-abstraction body, or the term
part of a type application. No reduction position lies inside a type.
This proves both one-step assertions.

For the last assertion, each free occurrence of \(x\) in the finite
original term receives a designated copy of \(N\). Contract the given
redex inside each such copy, using compatible reduction at that position.
The copies are disjoint; binders were renamed to prevent capture.
After finitely many such steps the result is \(M[N'/x]\). If no copy
was inserted the two terms coincide. This argument includes reductions
under both kinds of abstraction and within the term part of type
applications. Iterating the one-step facts gives their finite-step forms.

Finally a finite raw term has finitely many redex positions. Suppose it
has arbitrarily long reductions but no infinite reduction. Some immediate
successor must then have arbitrarily long reductions, because there are
only finitely many immediate successors. Repeat this choice, always taking
the first such successor in a fixed ordering of redex positions. It produces
an infinite reduction, a contradiction. Therefore a strongly normalizing
term has a finite greatest reduction length. For an immediate successor,
prefixing any of its reductions by the first step shows
\(h_{\mathsf F}(M')+1\leq h_{\mathsf F}(M)\).
\end{proof}
\end{annexitem}

\begin{annexitem}{annex:p7.system-f.typing-substitution}
 {thm:p07-system-f-sn-typing-substitution}{Formation, typing and preservation of the two redexes}
\begin{proof}
Type formation is preserved by substitution, by induction on the type:
a variable uses its specified image; an arrow uses the two induction
hypotheses; a universal type first renames its binder fresh for the target
context and substitution images and then uses the premise under that fresh
type variable. The same elementary inductions show type-context weakening,
term-context weakening and consistent renaming of fresh binders.

For the first typing assertion, induct on its typing derivation. A variable
uses its declaration with the substituted type in \(\Gamma\tau\).
Application uses the two substituted premise judgments and arrow elimination.
At \(\lambda x{:}\beta.P\), rename \(x\) fresh if necessary, apply
induction to the premise under \(x:\beta\tau\), and use arrow
introduction with annotation \(\beta\tau\).
At \(\Lambda a.P\), rename \(a\) fresh for \(\Delta'\), the
substitution domain and all substitution images. Extend \(\tau\) by
fixing \(a\). The premise induction works under \(\Delta',a\);
the context is \(\Gamma\tau\), since the source eigenvariable was
absent from every declared context type. The fresh \(a\) is also absent
from these substituted types, so universal introduction applies.
At \(P[\beta]\), the premise gives the substituted universal type,
type formation gives \(\beta\tau\), and universal elimination has
result type
\((\alpha\tau)[\beta\tau/a]=(\alpha[\beta/a])\tau\), after
renaming the universal binder fresh for \(\tau\). The syntactic
substitution induction in the preceding annex item proves this equality.

For the second assertion, again induct on the five typing cases. A variable
is precisely the hypothesis about its image. Application applies the two
induction hypotheses. A lambda abstraction first renames its binder \(x\)
fresh for \(\Gamma'\) and all images, weakens their typings to the context
\(\Gamma',x:\beta\), and extends the substitution by \(x\mapsto x\).
The premise induction followed by arrow introduction gives the required
substituted lambda. At type abstraction, choose its type binder \(a\)
fresh for the target context and the images. Type-context weakening types
the same images under \(\Delta,a\). The premise induction followed by
universal introduction is legitimate since \(a\) is absent from
the types in \(\Gamma'\). A type application retains its well-formed
type argument and uses the induction hypothesis for its function.
This proves term substitution with all freshness conditions.

At a typed term-beta redex, inversion of application and lambda typing
provides \(\Delta;\Gamma,x:\beta\derives_{\mathsf F}P:\alpha\)
and \(\Delta;\Gamma\derives_{\mathsf F}N:\beta\). Term substitution
therefore types \(P[N/x]\) at \(\alpha\) in \(\Gamma\).
At a typed type-beta redex, inversion provides
\(\Delta,a;\Gamma\derives_{\mathsf F}P:\alpha\), with \(a\) absent
from \(\Gamma\), and a type \(\beta\) over \(\Delta\). Type
substitution by \([\beta/a]\) gives
\(\Delta;\Gamma\derives_{\mathsf F}P[\beta/a]:\alpha[\beta/a]\).
These are exactly the types of the two redexes.
Induction on the surrounding typing derivation handles contractions in
either side of a term application, in a lambda body, in a type-abstraction
body, and in the term part of a type application. Each enclosing rule has
the same context, result type and side conditions after replacing its
affected premise. This proves subject reduction for the full relation.
\end{proof}
\end{annexitem}

\begin{annexitem}{annex:p7.system-f.candidate-constructors}
 {thm:p07-system-f-sn-candidate-constructors}{Closure properties of all three candidate constructors}
\begin{proof}
Write SN for \(\mathsf{SN}_{\mathsf F}\) in this proof. SN satisfies
clause~1 by definition. Any infinite reduction from a successor would
extend to one from its predecessor, proving clause~2. If all immediate
reducts of a neutral term are strongly normalizing, an infinite reduction
from that term would begin with one of those reducts, a contradiction.
This proves clause~3. For any candidate, a variable is neutral and has
no immediate reducts, so clause~3 includes it. Finite reduction closure
follows by induction on the number of steps.

Let \(R,S\) be candidates. If \(M\in R\Rightarrow S\), take any
variable \(z\in R\). Then \(Mz\in S\subseteq\mathrm{SN}\).
An infinite reduction of \(M\) would remain infinite in the function
position of \(Mz\); hence \(M\in\mathrm{SN}\), giving clause~1.
For clause~2, if \(M\to_{\mathsf F}M'\), then for any \(N\in R\)
we have \(MN\to_{\mathsf F}M'N\). Clause~2 of \(S\) gives
\(M'N\in S\), as required.

For clause~3, suppose \(M\) is neutral and every immediate reduct
\(M'\) belongs to \(R\Rightarrow S\). Fix \(N\in R\) and induct
on its finite height. Every immediate reduct of \(MN\) has one of two
forms. A reduction in the function gives \(M'N\in S\) by the
hypothesis on \(M'\). A reduction in the argument gives \(MN'\);
clause~2 of \(R\) gives \(N'\in R\), and its smaller height gives
\(MN'\in S\) by induction. There is no root term-beta contraction
because \(M\) is not a lambda abstraction; the whole term application
cannot be a root type-beta redex. Since \(MN\) is neutral, clause~3
of \(S\) now puts \(MN\) in \(S\). This proves the arrow clause.

For \(\mathcal U(H)\), take one fixed type \(b\) and the candidate
SN. Membership gives \(M[b]\in H_{b,\mathrm{SN}}\subseteq\mathrm{SN}\).
An infinite reduction in \(M\) would persist inside that type application,
so \(M\) is strongly normalizing. If \(M\to_{\mathsf F}M'\), then
\(M[\beta]\to_{\mathsf F}M'[\beta]\); closure of each
\(H_{\beta,C}\) gives clause~2. Finally suppose \(M\) is neutral
and every immediate reduct lies in \(\mathcal U(H)\). An immediate
reduct of \(M[\beta]\) is necessarily \(M'[\beta]\): types have
no reductions and neutrality excludes a root type-beta contraction.
Each such reduct belongs to \(H_{\beta,C}\). The whole type
application is neutral, so that candidate's clause~3 gives membership.
This holds for every \(\beta,C\), proving clause~3 and the result.
\end{proof}
\end{annexitem}

\begin{annexitem}{annex:p7.system-f.interpretation}
 {thm:p07-system-f-sn-interpretation}{Parametric interpretation and semantic substitution}
\begin{proof}
Structural induction on \(\alpha\) proves that its interpretation is
a candidate. A type variable uses the candidate supplied by its valuation.
An arrow uses the arrow constructor from the preceding lemma and the two
induction hypotheses. A universal type uses that lemma's universal
constructor with
\(H_{\beta,C}=\mathcal R_{\rho[a\mapsto C]}(\alpha)\).
For each index this is a candidate by induction on the proper subtype
\(\alpha\). The already defined set of all candidates, rather than an
interpretation of a potentially larger substituted type, is quantified over.

If two valuations agree at the free variables of a type, their
interpretations coincide. At a variable this is the agreement assumption;
at an arrow it follows on each component. At \(\forall a.\alpha\),
updating both valuations with the same candidate at \(a\) makes them
agree on every free variable of its body. The induction hypothesis gives
equal membership conditions for every type and candidate index.

For completeness, fresh renaming is also respected. If \(b\) does not
occur free in \(\alpha\), renaming the free \(a\)'s of \(\alpha\)
to \(b\), and interpreting the new \(b\) by \(C\), has the same
effect as interpreting the original \(a\) by \(C\). A structural
induction proves this: the variable \(a\) has value \(C\) on both
sides, other occurring variables retain their valuations, arrows use the
two component identities, and a universal binder is first renamed away
from \(a,b\) and then ranges over the same arbitrary candidate on both
sides. Apply this fact to the bodies of \(\forall a.\alpha\) and
\(\forall b.\alpha[b/a]\). Their universal membership conditions
coincide for each syntactic type argument and each candidate, proving
bound-renaming invariance.

Now prove the displayed substitution identity by induction on
\(\alpha\). At \(\alpha=a\), both sides are
\(\mathcal R_\rho(\beta)\). At any other variable the update does
not change its value. At an arrow apply the two induction hypotheses and
the arrow definition. In the universal case rename its binder \(b\)
fresh, so \(b\ne a\) and \(b\notin\operatorname{FV}(\beta)\).
An arbitrary raw term \(M\) belongs to the left-hand interpretation
exactly when, for every type \(\gamma\) and candidate \(C\),
\[
 M[\gamma]\in
 \mathcal R_{\rho[b\mapsto C]}(\delta[\beta/a]),
\]
where the original type is \(\forall b.\delta\). The induction
hypothesis turns the latter interpretation into
\[
 \mathcal R_{\rho[b\mapsto C]
     [a\mapsto\mathcal R_{\rho[b\mapsto C]}(\beta)]}(\delta).
\]
The free-variable agreement fact makes
\(\mathcal R_{\rho[b\mapsto C]}(\beta)=\mathcal R_\rho(\beta)\).
The two updates at distinct variables commute. Thus the membership
condition is precisely the universal interpretation under
\(\rho[a\mapsto\mathcal R_\rho(\beta)]\), proving the identity.
All candidate-valued parameters used in this argument exist by the first
part of the proof, including interpretations of impredicative instances.
\end{proof}
\end{annexitem}

\begin{annexitem}{annex:p7.system-f.beta-expansion}
 {thm:p07-system-f-sn-beta-expansion}{Exhausting the immediate reducts of both beta forms}
\begin{proof}
For the term-beta form induct on
\(h_{\mathsf F}(P)+h_{\mathsf F}(N)\). Its root contractum is
\(P[N/x]\in R\). A step in \(P\) gives
\((\lambda x{:}\alpha.P')N\). Substitution preservation gives
\(P[N/x]\to_{\mathsf F}P'[N/x]\), so candidate closure puts the
new contractum in \(R\); the smaller height permits induction.
A step in \(N\) gives \((\lambda x{:}\alpha.P)N'\).
The finite-copy substitution lemma gives
\(P[N/x]\to_{\mathsf F}^*P[N'/x]\), possibly in zero steps.
Candidate closure and the smaller height again permit induction.
There is no reduction in the annotation \(\alpha\), and these are
all immediate reducts of the redex. The whole application is neutral;
clause~3 therefore places it in \(R\).

For the type-beta form induct on \(h_{\mathsf F}(P)\). The root
contractum is \(P[\beta/a]\in R\). The only other possibility is
\((\Lambda a.P')[\beta]\) with \(P\to_{\mathsf F}P'\).
Type substitution preserves that step, and clause~2 gives
\(P'[\beta/a]\in R\). The smaller height permits induction.
There is no reduction in \(\beta\). Every immediate reduct of this
neutral type application therefore belongs to \(R\), so clause~3
applies. These arguments concern the combined reduction in \(P\) and
\(N\), not just one of its two generating beta rules.
\end{proof}
\end{annexitem}

\begin{annexitem}{annex:p7.system-f.introduction}
 {thm:p07-system-f-sn-introduction}{Recovering the normalization premises for introduction}
\begin{proof}
For arrow introduction, choose a term variable \(z\) fresh for \(P\).
It belongs to \(R\). The hypothesis therefore puts \(P[z/x]\) in
\(S\subseteq\mathsf{SN}_{\mathsf F}\). An infinite reduction of
\(P\) would map to an infinite sequence of genuine contractions in
\(P[z/x]\), by substitution preservation. Thus \(P\) is strongly
normalizing. Every \(N\in R\) is also strongly normalizing; controlled
term-beta expansion now puts \((\lambda x{:}\alpha.P)N\) in \(S\)
for each such \(N\). This is the arrow membership condition.

For universal introduction, first take the syntactic type argument to be
the variable \(a\) itself and take the candidate to be
\(\mathsf{SN}_{\mathsf F}\). The hypothesis gives
\[
 P=P[a/a]\in
 \mathcal R_{\rho[a\mapsto\mathsf{SN}_{\mathsf F}]}(\alpha)
 \subseteq\mathsf{SN}_{\mathsf F}.
\]
This justifies the normalization premise rather than assuming it.
For any type \(\beta\) and candidate \(C\), the given contractum
\(P[\beta/a]\) belongs to the candidate
\(\mathcal R_{\rho[a\mapsto C]}(\alpha)\). Controlled type-beta
expansion puts \((\Lambda a.P)[\beta]\) in that same candidate.
Quantifying over all \(\beta,C\) is exactly the universal
interpretation's membership condition.
\end{proof}
\end{annexitem}

\begin{annexitem}{annex:p7.system-f.fundamental}
 {thm:p07-system-f-sn-fundamental}{All five typing cases with both substitutions}
\begin{proof}
Induct on the typing derivation, with the stated conclusion universally
quantified over the type substitution, valuation and candidate-respecting
term substitution. All variable sets used in a particular derivation and
its substitution images are finite, so fresh binders exist.

\paragraph{Variable.}
If the last rule reads \(x:\alpha\) from the context, then
\((x\tau)\sigma=\sigma(x)\in\mathcal R_\rho(\alpha)\), precisely
the substitution hypothesis.

\paragraph{Term application.}
The two premises type \(P:\beta\mathbin{\to_{\rm ty}}\alpha\)
and \(N:\beta\). Induction puts their doubly substituted terms in
the corresponding arrow and argument candidates. By the definition of
\(\Rightarrow\), their application is in \(\mathcal R_\rho(\alpha)\).
Both substitutions distribute over application, so this is the desired
term.

\paragraph{Term abstraction.}
Write the term as \(\lambda x{:}\beta.P\), with premise typing
\(P:\gamma\) under \(\Gamma,x:\beta\). Rename \(x\) fresh for
the term-substitution domain and the free term variables of all images.
For every \(N\in\mathcal R_\rho(\beta)\), extend \(\sigma\) by
\(x\mapsto N\). Induction on the premise gives
\[
 (P\tau)(\sigma[x\mapsto N])\in\mathcal R_\rho(\gamma).
\]
By the freshness condition, the left side is
\(((P\tau)\sigma)[N/x]\). Arrow introduction in the interpretation
therefore gives
\[
 \lambda x{:}\beta\tau.((P\tau)\sigma)
 \in\mathcal R_\rho(\beta)\Rightarrow\mathcal R_\rho(\gamma).
\]
This is the doubly substituted abstraction, with its actual substituted
annotation, as required. If an argument contains \(x\) freely, its
image is still inserted once and capture-avoiding substitution gives the
same identity; closed arguments were not assumed.

\paragraph{Type abstraction.}
Write the term as \(\Lambda a.P\), with premise
\(\Delta,a;\Gamma\derives_{\mathsf F}P:\beta\).
The sheet requires \(a\) absent from \(\Delta\) and from every
type in \(\Gamma\). Rename it further to be fresh for all free type
variables in the images of \(\tau\) and \(\sigma\). The valuation
interprets bound type names only under its local updates, so this renaming
does not affect the conclusion.

Fix an arbitrary syntactic type \(\gamma\) and arbitrary candidate
\(C\). Extend the type substitution by \(a\mapsto\gamma\) and
change the valuation to \(\rho[a\mapsto C]\). For every
\(x:\delta\in\Gamma\), the source side condition and the semantic
free-variable lemma give
\(\mathcal R_{\rho[a\mapsto C]}(\delta)=\mathcal R_\rho(\delta)\).
Thus the same \(\sigma\) satisfies the premise induction hypothesis,
which gives
\[
 (P(\tau[a\mapsto\gamma]))\sigma
     \in\mathcal R_{\rho[a\mapsto C]}(\beta).
\]
Freshness for the images ensures the syntactic equality
\[
 (P(\tau[a\mapsto\gamma]))\sigma
       =((P\tau)\sigma)[\gamma/a].
\]
Indeed type substitution commutes with the other substitutions after the
fresh renaming: no inserted type or term image contains a free \(a\)
that the final substitution could otherwise change. The equality follows
by the five-constructor substitution induction used earlier.
Universal introduction in the interpretation now yields
\(\Lambda a.((P\tau)\sigma)\in\mathcal R_\rho(\forall a.\beta)\).
This is the required substituted term. No restriction was imposed on the
complexity or free variables of the instance \(\gamma\).

\paragraph{Type application.}
The last rule types \(P[\gamma]\) from
\(P:\forall a.\beta\), with result \(\beta[\gamma/a]\).
Induction puts \((P\tau)\sigma\) in the universal interpretation.
Choose its syntactic argument to be \(\gamma\tau\) and its
candidate argument to be \(\mathcal R_\rho(\gamma)\), which is a
candidate by the interpretation lemma. This gives
\[
 ((P\tau)\sigma)[\gamma\tau]
 \in\mathcal R_{\rho[a\mapsto\mathcal R_\rho(\gamma)]}(\beta)
 =\mathcal R_\rho(\beta[\gamma/a]).
\]
The equality is semantic type substitution. The term on the left is
\(((P[\gamma])\tau)\sigma\), since term substitution leaves a
type argument unchanged. This covers universal elimination for all
well-formed instances, including impredicative ones, and completes the
induction over every typing rule.
\end{proof}
\end{annexitem}

\begin{annexitem}{annex:p7.system-f.strong-normalization}
 {thm:p07-system-f-strong-normalization}{Strong normalization of the exact open Church-style calculus}
\begin{proof}
Take any derivation \(\Delta;\Gamma\derives_{\mathsf F}M:\alpha\)
on Sheet~\ref{sheet:system-f}. Assign the candidate
\(\mathsf{SN}_{\mathsf F}\) to every free type variable, choose the
identity type substitution on \(\Delta\), and choose the identity
term substitution on the variables of \(\Gamma\). Every variable is
in every candidate, hence in the interpretation of its assigned type.
The fundamental lemma therefore gives
\(M\in\mathcal R_\rho(\alpha)\). This interpretation is a candidate,
so its first clause gives \(M\in\mathsf{SN}_{\mathsf F}\).
By definition there is no infinite sequence of either kind of beta
contraction, interleaved in any order and at any term positions.

Neither substitution closed the free variables, so the argument applies
to open judgments exactly as stated. A maximal reduction sequence must
be finite; if its last term had a reduct the sequence could be extended,
contradicting maximality. It therefore ends in a normal term. Iterating
subject reduction preserves the type and both contexts along this sequence.
No confluence premise, erased-term termination premise or predicative
restriction is used to obtain termination.
\end{proof}
\end{annexitem}
